%==============================================================================
%  UG PROJECT REPORT  --  Chandigarh University format
%  Project : An Interpretable Multi-Modal Spam Detection Framework Using
%            Explainable AI and Machine Learning
%  Compile with: pdfLaTeX  (Overleaf: Menu -> Compiler -> pdfLaTeX)
%  Compile 2-3 times so the Table of Contents / List of Figures fill in.
%
%  IMAGES live in the figures/ sub-folder (on Overleaf, upload the folder as is).
%  The document compiles even if one is missing (a labelled grey box is shown):
%     figures/cu_logo.png              <- Chandigarh University logo (cover+certificate)
%     figures/dataset_overview.png     <- class balance & message length (EDA)
%     figures/sms_roc_curves.png       <- ROC curves, SMS models
%     figures/email_roc_curves.png     <- ROC curves, e-mail models
%     figures/sms_shap_waterfall.png   <- SHAP waterfall of the exhibited SMS
%     figures/sms_shap_beeswarm.png    <- global SHAP beeswarm, SMS
%     figures/email_shap_beeswarm.png  <- global SHAP beeswarm, e-mail
%     figures/app_detector_tab.png     <- screenshot of the Streamlit Detector tab
%  The three SHAP figures are generated by scripts/make_shap_figures.py.
%==============================================================================
\documentclass[12pt,a4paper]{report}

%------------------------- page geometry (1 inch all round) -------------------
\usepackage[a4paper,top=1in,bottom=1in,left=1in,right=1in]{geometry}
% For a binding margin instead, use:
% \usepackage[a4paper,top=1in,bottom=1in,left=1.25in,right=1in]{geometry}

%------------------------- fonts: Times New Roman clone -----------------------
\usepackage[T1]{fontenc}
\usepackage[utf8]{inputenc}
\usepackage{mathptmx}          % Nimbus Roman = metrically identical to Times New Roman

%------------------------- spacing, layout, tools -----------------------------
\usepackage{setspace}
\usepackage{graphicx}
\usepackage{xcolor}
\usepackage{titlesec}
\usepackage{tocloft}
\usepackage[font={bf,footnotesize},labelsep=colon,justification=centering,skip=6pt]{caption}
\usepackage{enumitem}
\usepackage{array}
\usepackage{longtable}
\usepackage{float}
\usepackage{multirow}
\usepackage{amsmath,amssymb}
\usepackage{pgfplots}                 % charts drawn directly in LaTeX
\pgfplotsset{compat=1.17}
\usepackage{pgfgantt}                 % Gantt chart for the timeline
\usetikzlibrary{arrows.meta,positioning,calc,shapes.geometric}
\usepackage[hidelinks,bookmarksnumbered]{hyperref}   % keep hyperref last

\onehalfspacing                        % 1.5 line spacing for the general text
\setlength{\parindent}{0pt}            % Word-style: no first-line indent
\setlength{\parskip}{6pt}
\setlist{nosep,topsep=4pt,itemsep=3pt,leftmargin=1.2cm}
\renewcommand{\arraystretch}{1.3}
% ragged-right paragraph column for tables:  L{3cm}  instead of  p{3cm}
\newcolumntype{L}[1]{>{\raggedright\arraybackslash}p{#1}}

% Word does not hyphenate by default -- switch it off so the justification matches
\frenchspacing          % Word-style: no extra space after a full stop
\hyphenpenalty=10000
\exhyphenpenalty=10000
\tolerance=9999
\emergencystretch=3em

%------------------------- exact font sizes from the format sheet -------------
\newcommand{\sixteen}{\fontsize{16}{19.2}\selectfont}   % chapter names
\newcommand{\fourteen}{\fontsize{14}{21}\selectfont}    % main headings
\newcommand{\twelve}{\fontsize{12}{18}\selectfont}      % body / sub-headings
\newcommand{\eighteen}{\fontsize{18}{27}\selectfont}    % report title

%------------------------- CHAPTER: 16pt, bold, centered ----------------------
\titleformat{\chapter}[display]
  {\filcenter\bfseries\sixteen}
  {CHAPTER \thechapter.}
  {6pt}
  {}
\titlespacing*{\chapter}{0pt}{-24pt}{22pt}

% unnumbered chapters (REFERENCES, APPENDIX, USER MANUAL ...)
\titleformat{name=\chapter,numberless}[display]
  {\filcenter\bfseries\sixteen}{}{0pt}{}
\titlespacing*{name=\chapter,numberless}{0pt}{-24pt}{22pt}

%------------------------- HEADING: 14pt, bold, left ---------------------------
\titleformat{\section}[hang]
  {\bfseries\fourteen}
  {\makebox[3.4em][l]{\thesection.}}{0pt}{}
\titlespacing*{\section}{0pt}{16pt}{8pt}

%------------------------- SUB-HEADING: 12pt, bold, left -----------------------
\titleformat{\subsection}[hang]
  {\bfseries\twelve}
  {\makebox[3.4em][l]{\thesubsection.}}{0pt}{}
\titlespacing*{\subsection}{0pt}{12pt}{6pt}

\titleformat{\subsubsection}[hang]
  {\bfseries\twelve}
  {\makebox[3.9em][l]{\thesubsubsection.}}{0pt}{}
\titlespacing*{\subsubsection}{0pt}{10pt}{6pt}

\setcounter{secnumdepth}{3}
\setcounter{tocdepth}{2}

%------------------------- TABLE OF CONTENTS styling ---------------------------
\renewcommand{\contentsname}{TABLE OF CONTENTS}
\renewcommand{\cfttoctitlefont}{\bfseries\sixteen}
\setlength{\cftbeforetoctitleskip}{0pt}
\setlength{\cftaftertoctitleskip}{22pt}

% "Chapter 1. ....... 4"
\renewcommand{\cftchappresnum}{Chapter\ }
\renewcommand{\cftchapaftersnum}{.}
\setlength{\cftchapnumwidth}{6.2em}
\renewcommand{\cftchapfont}{\bfseries}
\renewcommand{\cftchappagefont}{\bfseries}
\renewcommand{\cftchapleader}{\bfseries\cftdotfill{\cftdotsep}}
\renewcommand{\cftchapdotsep}{\cftdotsep}
\setlength{\cftbeforechapskip}{8pt}

\renewcommand{\cftsecfont}{\bfseries}
\renewcommand{\cftsecpagefont}{\bfseries}
\renewcommand{\cftsecleader}{\bfseries\cftdotfill{\cftdotsep}}
\setlength{\cftsecnumwidth}{3.2em}
\setlength{\cftbeforesecskip}{6pt}

\renewcommand{\cftsubsecfont}{\bfseries}
\renewcommand{\cftsubsecleader}{\bfseries\cftdotfill{\cftdotsep}}
\renewcommand{\cftsubsecpagefont}{\bfseries}
\setlength{\cftsubsecnumwidth}{4.0em}
\setlength{\cftbeforesubsecskip}{6pt}

%------------------------- LIST OF FIGURES / TABLES styling --------------------
\renewcommand{\listfigurename}{List of Figures}
\renewcommand{\cftloftitlefont}{\bfseries\sixteen}
\setlength{\cftbeforeloftitleskip}{0pt}
\setlength{\cftafterloftitleskip}{22pt}
\renewcommand{\cftfigpresnum}{Figure\ }
\renewcommand{\cftfigaftersnum}{:}
\setlength{\cftfignumwidth}{5.6em}
\renewcommand{\cftfigfont}{\bfseries}
\renewcommand{\cftfigpagefont}{\bfseries}
\renewcommand{\cftfigleader}{\bfseries\cftdotfill{\cftdotsep}}
\setlength{\cftbeforefigskip}{8pt}

\renewcommand{\listtablename}{List of Tables}
\renewcommand{\cftlottitlefont}{\bfseries\sixteen}
\setlength{\cftbeforelottitleskip}{0pt}
\setlength{\cftafterlottitleskip}{22pt}
\renewcommand{\cfttabpresnum}{Table\ }
\renewcommand{\cfttabaftersnum}{:}
\setlength{\cfttabnumwidth}{5.2em}
\renewcommand{\cfttabfont}{\bfseries}
\renewcommand{\cfttabpagefont}{\bfseries}
\renewcommand{\cfttableader}{\bfseries\cftdotfill{\cftdotsep}}
\setlength{\cftbeforetabskip}{8pt}

% centre the three list titles (tocloft's own \hfill trick loses to \parfillskip)
\makeatletter
\renewcommand{\@cftmaketoctitle}{%
  \addpenalty\@secpenalty \vspace*{\cftbeforetoctitleskip}\@cftpagestyle
  {\interlinepenalty\@M {\centering\cfttoctitlefont\contentsname\par}%
   \cftmarktoc \nobreak \vskip\cftaftertoctitleskip \@afterheading}}
\renewcommand{\@cftmakeloftitle}{%
  \addpenalty\@secpenalty \vspace*{\cftbeforeloftitleskip}\@cftpagestyle
  {\interlinepenalty\@M {\centering\cftloftitlefont\listfigurename\par}%
   \cftmarklof \nobreak \vskip\cftafterloftitleskip \@afterheading}}
\renewcommand{\@cftmakelottitle}{%
  \addpenalty\@secpenalty \vspace*{\cftbeforelottitleskip}\@cftpagestyle
  {\interlinepenalty\@M {\centering\cftlottitlefont\listtablename\par}%
   \cftmarklot \nobreak \vskip\cftafterlottitleskip \@afterheading}}
\makeatother

%------------------------- helper: logo with safe fallback ---------------------
\newcommand{\culogo}[1][2.6in]{%
  \IfFileExists{figures/cu_logo.png}
    {\includegraphics[width=#1]{figures/cu_logo.png}}
    {\fbox{\parbox[c][1.0in][c]{#1}{\centering\small\itshape
      upload \texttt{figures/cu\_logo.png}\\to the project folder}}}%
}

%------------------------- helper: image with safe fallback --------------------
% \figimg[width]{filename}{fallback-box-height}
\newcommand{\figimg}[3][0.75\textwidth]{%
  \IfFileExists{#2}
    {\includegraphics[width=#1]{#2}}
    {\fbox{\parbox[c][#3][c]{#1}{\centering\small\itshape
      upload \texttt{\detokenize{#2}}\\to the project folder}}}%
}

%------------------------- TikZ styles for the diagrams -------------------------
\tikzset{
  blk/.style={draw=black!70, fill=gray!8, rounded corners=2pt, align=center,
              font=\small, minimum height=0.85cm, inner sep=5pt},
  mdl/.style={blk, fill=blue!8},
  xai/.style={blk, fill=orange!14},
  dat/.style={blk, fill=green!10},
  arr/.style={-{Latex[length=2.6mm]}, semithick},
}

%------------------------- front-matter page helper ----------------------------
\newcommand{\frontchapter}[1]{%
  \chapter*{#1}\phantomsection\addcontentsline{toc}{chapter}{#1}%
}

%==============================================================================
%  >>>>>>>>>>>>>>>>>>>>>  EDIT YOUR DETAILS HERE  <<<<<<<<<<<<<<<<<<<<<<<<<<<<
%==============================================================================
\newcommand{\ProjectTitle}{AN INTERPRETABLE MULTI-MODAL SPAM DETECTION FRAMEWORK USING EXPLAINABLE AI AND MACHINE LEARNING}
\newcommand{\CandidateA}{Aariz (23BCS10422)}
\newcommand{\CandidateB}{Naman (23BCS11036)}
\newcommand{\CandidateC}{Vanshit Kumar (23BCS10087)}
\newcommand{\CandidateD}{Aditi Arya (23BCS11535)}
\newcommand{\CandidateE}{Avish (23BCS13898)}
\newcommand{\AllCandidates}{Aariz (23BCS10422), Naman (23BCS11036), Vanshit Kumar (23BCS10087), Aditi Arya (23BCS11535) and Avish (23BCS13898)}
\newcommand{\DegreeName}{BACHELOR OF ENGINEERING}
\newcommand{\BranchName}{COMPUTER SCIENCE AND ENGINEERING}
\newcommand{\MonthYear}{NOVEMBER 2026}                      % <-- EDIT month of submission
\newcommand{\HoDName}{Dr.~Surinder Kaur}                % <-- EDIT: fill in your HoD's name
\newcommand{\SupervisorName}{Sonam Juneja (E11002)}
\newcommand{\SupervisorDesignation}{Assistant Professor}    % <-- verify designation
\newcommand{\DeptName}{Department of Computer Science \& Engineering}
%==============================================================================

\begin{document}

%==============================================================================
%  COVER / TITLE PAGE
%==============================================================================
\pagestyle{empty}
\pagenumbering{roman}

\begingroup
\setlength{\parskip}{0pt}
\singlespacing
\begin{center}
\vspace*{0.12in}

{\bfseries\eighteen \ProjectTitle\par}

\vspace{0.34in}
{\bfseries\fourteen A PROJECT REPORT\par}

\vspace{0.30in}
{\bfseries\itshape\fourteen Submitted by\par}

\vspace{0.26in}
{\bfseries\sixteen \CandidateA\par}
\vspace{0.06in}
{\bfseries\sixteen \CandidateB\par}
\vspace{0.06in}
{\bfseries\sixteen \CandidateC\par}
\vspace{0.06in}
{\bfseries\sixteen \CandidateD\par}
\vspace{0.06in}
{\bfseries\sixteen \CandidateE\par}

\vspace{0.30in}
{\bfseries\itshape\fourteen in partial fulfillment for the award of the degree of\par}

\vspace{0.28in}
{\bfseries\sixteen \DegreeName\par}

\vspace{0.16in}
{\bfseries\fourteen IN\par}

\vspace{0.16in}
{\fourteen \BranchName\par}

\vspace{0.28in}
\culogo

\vspace{0.10in}
{\bfseries\fourteen Chandigarh University\par}

\vspace{0.22in}
{\fourteen \MonthYear\par}

\end{center}
\endgroup
\clearpage

%==============================================================================
%  BONAFIDE CERTIFICATE
%==============================================================================
\begin{center}
\vspace*{0.1in}
\culogo

\vspace{0.35in}
{\bfseries\sixteen BONAFIDE CERTIFICATE\par}
\end{center}

\vspace{0.35in}

{\fourteen
Certified that this project report \textbf{``\ProjectTitle''} is the bonafide
work of \textbf{``\AllCandidates''} who carried out the project work under
my/our supervision.\par}

\vspace{0.45in}

\begin{singlespace}
\noindent
\begin{minipage}[t]{0.47\textwidth}
\raggedright
\fourteen
\vspace*{0.45in}
\textbf{SIGNATURE}\\[0.18in]
\HoDName\\[0.10in]
\textbf{HEAD OF THE DEPARTMENT}\\[0.10in]
\DeptName
\end{minipage}
\hfill
\begin{minipage}[t]{0.47\textwidth}
\raggedright
\fourteen
\vspace*{0.45in}
\textbf{SIGNATURE}\\[0.18in]
\SupervisorName\\[0.10in]
\textbf{SUPERVISOR}\\[0.10in]
\SupervisorDesignation\\[0.10in]
\DeptName
\end{minipage}
\end{singlespace}

\vspace{0.45in}

{\fourteen Submitted for the project viva-voce examination held on
\rule{1.6in}{0.4pt}\par}

\vspace{0.45in}

\begin{singlespace}
\noindent
\begin{minipage}[t]{0.47\textwidth}
\raggedright\bfseries\fourteen INTERNAL EXAMINER
\end{minipage}
\hfill
\begin{minipage}[t]{0.47\textwidth}
\raggedright\bfseries\fourteen EXTERNAL EXAMINER
\end{minipage}
\end{singlespace}

\clearpage

%==============================================================================
%  ACKNOWLEDGEMENT
%==============================================================================
\begin{center}
\vspace*{0.1in}
{\bfseries\sixteen ACKNOWLEDGEMENT\par}
\end{center}
\vspace{0.3in}

We would like to express our sincere gratitude to our supervisor,
\SupervisorName, \SupervisorDesignation, \DeptName, Chandigarh University, for
the continuous guidance, constructive criticism and encouragement extended
throughout the course of this project. Her insistence that every design
decision be justified --- and every prediction be explainable --- shaped the
central theme of this work.

We are equally thankful to \HoDName, Head of the Department, for providing the
computing facilities and the academic environment that made this work
possible. Our thanks also go to the faculty members of the department for
their timely feedback during the project reviews.

Finally, we thank our families and friends for their constant support and
motivation.

\vspace{0.6in}
\hfill\begin{minipage}{0.5\textwidth}
\raggedright
\CandidateA\\
\CandidateB\\
\CandidateC\\
\CandidateD\\
\CandidateE
\end{minipage}

\clearpage

%==============================================================================
%  TABLE OF CONTENTS
%==============================================================================
\tableofcontents
\addtocontents{toc}{\protect\thispagestyle{empty}}
\clearpage

%==============================================================================
%  From here the roman page numbers become visible: i, ii, iii ...
%==============================================================================
\pagestyle{plain}
\setcounter{page}{1}

\phantomsection\addcontentsline{toc}{chapter}{List of Figures}
\listoffigures
\clearpage

\phantomsection\addcontentsline{toc}{chapter}{List of Tables}
\listoftables
\clearpage

%==============================================================================
%  ABSTRACT
%==============================================================================
\frontchapter{ABSTRACT}

Fraudulent messages now reach users through more than one channel, and the
machine-learning filters that intercept them are usually black boxes: they
return a verdict but no reason. Aggregate accuracy alone can conceal
systematic errors and dataset artefacts, leaves the end user with no basis to
contest a decision, and gives an analyst nothing to audit.

This project presents an interpretable multi-modal spam detection framework
covering both SMS and e-mail. Each modality is served by its own specialist
pipeline: text normalisation that maps high-signal entities to readable
placeholder tokens (\texttt{httpaddr}, \texttt{numbr}, \texttt{moneysymb},
\texttt{phonenumbr}), TF--IDF features over uni- and bi-grams combined with
fifteen engineered stylometric statistics, and a hybrid ensemble of Logistic
Regression, Random Forest and XGBoost fused by performance-weighted soft
voting. A lightweight router selects the appropriate specialist at inference
time. A dual explanation layer justifies every prediction: LIME explains the
complete ensemble locally, while exact TreeSHAP provides local and global
attributions for the boosted-tree component. The framework ships as a fully
reproducible Python pipeline and a four-tab interactive Streamlit application.

Evaluated on 17,158 messages from two public research corpora using an 80/20
stratified split and 5-fold cross-validation, the hybrid ensemble achieved
98.55\% accuracy (ROC-AUC 0.9941) on the UCI SMS Spam Collection and 98.25\%
accuracy (ROC-AUC 0.9983) on the Enron-Spam corpus, matching (to four
decimal places) or beating every base learner on ROC-AUC. Beyond the numbers, the explanation layer exposed a
genuine dataset bias --- the token \texttt{enron} acting as a shortcut for
legitimate mail --- that no accuracy table could reveal, demonstrating the
practical value of shipping explanations together with predictions.

\vspace{12pt}
\textbf{Keywords:} spam detection, explainable AI, ensemble learning, SHAP,
LIME, TF--IDF, multi-modal classification, machine learning.

\clearpage

%==============================================================================
%  GRAPHICAL ABSTRACT
%==============================================================================
\frontchapter{GRAPHICAL ABSTRACT}

\vspace{0.5in}
\begin{center}
\begin{tikzpicture}[node distance=0.55cm and 0.85cm]
  \node[dat, minimum width=2.6cm] (msg) {SMS / E-mail\\message};
  \node[blk, right=of msg, minimum width=2.8cm] (feat)
       {TF--IDF (1--2 grams)\\+ 15 stylometric\\features};
  \node[mdl, right=of feat, minimum width=3.1cm] (ens)
       {Hybrid ensemble\\LR + RF + XGBoost\\(weighted soft vote)};
  \node[blk, right=of ens, minimum width=2.6cm, fill=green!12] (out)
       {Verdict +\\spam probability};
  \node[xai, below=of ens, minimum width=3.1cm] (xai)
       {LIME + SHAP\\``\emph{why} was this\\flagged?''};
  \draw[arr] (msg) -- (feat);
  \draw[arr] (feat) -- (ens);
  \draw[arr] (ens) -- (out);
  \draw[arr] (ens) -- (xai);
  \draw[arr] (feat.south) |- (xai.west);
\end{tikzpicture}
\end{center}

\vspace{0.35in}
Every message is routed to a modality-specific feature space, scored by a
hybrid ensemble of three heterogeneous learners, and returned to the user
together with a human-readable explanation of the verdict.

\clearpage

%==============================================================================
%  ABBREVIATIONS
%==============================================================================
\frontchapter{ABBREVIATIONS}

\begin{singlespace}
\begin{center}
\begin{tabular}{@{}p{3.5cm}p{9.5cm}@{}}
\textbf{AI}      & Artificial Intelligence \\
\textbf{AUC}     & Area Under the (ROC) Curve \\
\textbf{CV}      & Cross-Validation \\
\textbf{EDA}     & Exploratory Data Analysis \\
\textbf{FN}      & False Negative \\
\textbf{FP}      & False Positive \\
\textbf{LIME}    & Local Interpretable Model-agnostic Explanations \\
\textbf{LR}      & Logistic Regression \\
\textbf{ML}      & Machine Learning \\
\textbf{NLP}     & Natural Language Processing \\
\textbf{RF}      & Random Forest \\
\textbf{ROC}     & Receiver Operating Characteristic \\
\textbf{SHAP}    & SHapley Additive exPlanations \\
\textbf{SMS}     & Short Message Service \\
\textbf{TF--IDF} & Term Frequency--Inverse Document Frequency \\
\textbf{TN}      & True Negative \\
\textbf{TP}      & True Positive \\
\textbf{UCI}     & University of California, Irvine (ML Repository) \\
\textbf{XAI}     & Explainable Artificial Intelligence \\
\textbf{XGB}     & eXtreme Gradient Boosting (XGBoost) \\
\end{tabular}
\end{center}
\end{singlespace}

\clearpage

%==============================================================================
%  SYMBOLS
%==============================================================================
\frontchapter{SYMBOLS}

\begin{singlespace}
\begin{center}
\begin{tabular}{@{}p{3.5cm}p{9.5cm}@{}}
$x$                          & Input message (raw text) \\
$\hat{y}$                    & Predicted class label (spam / ham) \\
$p_m(\mathrm{spam}\mid x)$   & Spam probability output by base learner $m$ \\
$w_m$                        & Ensemble weight of base learner $m$ \\
$w_{t,d}$                    & TF--IDF weight of term $t$ in document $d$ \\
$P$                          & Precision \\
$R$                          & Recall \\
$F1$                         & Harmonic mean of precision and recall \\
$\phi_i$                     & SHAP attribution of feature $i$ \\
$N$                          & Number of documents in the corpus \\
\end{tabular}
\end{center}
\end{singlespace}

\clearpage

%==============================================================================
%  MAIN TEXT
%==============================================================================
\pagenumbering{arabic}
\setcounter{page}{1}

%------------------------------------------------------------------------------
\chapter{INTRODUCTION}

\section{Client Identification/Need Identification/Identification of relevant Contemporary issue}

Unsolicited and fraudulent messaging is a documented, growing contemporary
problem. The FBI's Internet Crime Complaint Center recorded close to
3,00,000 phishing/spoofing complaints in calendar year 2023 --- the single
largest crime category in its annual report --- and industry telemetry has for
years placed unsolicited traffic at roughly half of all e-mail exchanged
worldwide. The problem is equally visible on the SMS channel: as banking,
commerce and one-time-password flows moved to text messaging, so did the
scams (``smishing''). In India, the volume of unsolicited commercial
communication led the Telecom Regulatory Authority of India to notify the
Telecom Commercial Communications Customer Preference Regulations, 2018, and
to mandate a distributed-ledger registration system for all bulk senders ---
regulatory evidence that the issue is both real and unresolved.

Filters exist, but a second, newer need has emerged around \emph{transparency}.
Modern spam classifiers are statistical black boxes: a message disappears and
nobody --- not the end user, not the platform analyst, not the auditor --- can
say why. Transparency obligations in data-protection law (Article~22 of the
GDPR in Europe; the Digital Personal Data Protection Act, 2023 in India) are
pushing automated decision systems toward explainability, and a filter that
cannot justify its verdicts is increasingly difficult to defend.

The client for this work is any operator of a user-facing messaging channel
--- typified by a university IT cell that receives phishing-report tickets
from students and staff --- who needs a filter that is (a)~accurate across
more than one message channel, (b)~able to attach a human-readable
justification to every verdict so that flagged messages can be triaged and
contested, and (c)~cheap enough to train and audit on commodity hardware.
This is a consultancy-style need: the client does not merely want a spam
score, they want a \emph{defensible} spam score.

\section{Identification of Problem}

No openly available framework simultaneously (i)~handles both SMS and e-mail,
whose statistics differ radically (median message length of 12 words versus
155; spam prevalence of about 12\% versus about 48\% in the public
research corpora), (ii)~matches the state-of-the-art accuracy of classical
machine-learning filters, and (iii)~attaches a faithful, per-prediction
explanation to every verdict. The consequences of this gap are concrete:
verdicts cannot be verified by non-experts, misclassified legitimate messages
cannot be contested, and hidden dataset biases pass silently into deployed
models. The problem addressed by this project is therefore the design,
implementation and validation of an accurate, fully interpretable,
multi-modal spam detection framework.

\section{Identification of Tasks}

The work was decomposed into the following tasks:

\begin{enumerate}
  \item Survey of the literature on spam/scam detection and on post-hoc
        explainable-AI techniques.
  \item Acquisition, cleaning and de-duplication of two public corpora (UCI
        SMS Spam Collection; Enron-Spam), followed by exploratory data
        analysis.
  \item Design of an interpretability-preserving preprocessing scheme and
        feature space (readable placeholder tokens; engineered stylometric
        features).
  \item Training and tuning of three heterogeneous base learners per modality
        with class-imbalance handling.
  \item Design and implementation of the performance-weighted soft-voting
        ensemble and the modality router.
  \item Integration of the explanation layer (LIME for the full ensemble;
        exact TreeSHAP for the boosted-tree component) at local and global
        scope.
  \item Development of the four-tab interactive Streamlit demonstration
        application and an automated end-to-end smoke-test suite.
  \item Evaluation on held-out test sets, error and bias analysis, and
        documentation of the results in this report.
\end{enumerate}

\section{Timeline}

The project was executed over one semester of approximately fifteen working
weeks. The schedule is represented as a Gantt chart in
Figure~\ref{fig:gantt}; the model-training, explanation-layer and application
workstreams deliberately overlap between weeks 8 and 12 because the
explanation artefacts were used to debug the models as they were being
trained.

\begin{figure}[H]
  \centering
  \resizebox{0.98\textwidth}{!}{%
  \begin{ganttchart}[
      x unit=0.62cm, y unit title=0.62cm, y unit chart=0.56cm,
      canvas/.append style={draw=black!60},
      title/.append style={draw=black!60, fill=gray!12},
      title label font=\scriptsize\bfseries,
      bar/.append style={fill=gray!35, draw=black!70},
      bar height=0.55,
      bar label font=\scriptsize
    ]{1}{15}
    \gantttitle{Project week}{15}\\
    \gantttitlelist{1,...,15}{1}\\
    \ganttbar{Literature survey}{1}{3}\\
    \ganttbar{Data acquisition \& EDA}{2}{4}\\
    \ganttbar{Preprocessing \& feature design}{4}{6}\\
    \ganttbar{Model training \& ensembling}{6}{9}\\
    \ganttbar{XAI layer (SHAP + LIME)}{8}{11}\\
    \ganttbar{Streamlit demo application}{10}{12}\\
    \ganttbar{Validation \& smoke testing}{12}{14}\\
    \ganttbar{Report writing}{13}{15}
  \end{ganttchart}}
  \caption{Gantt chart of the project schedule}
  \label{fig:gantt}
\end{figure}

\section{Organization of the Report}

Chapter~2 reviews the reported literature on spam detection and explainable
AI and closes with the formal problem definition and objectives. Chapter~3
covers the evaluation of candidate features, the design constraints, the two
alternative design flows that were considered and the justification for the
selected design. Chapter~4 presents the implementation using modern
engineering tools together with the complete experimental results, error
analysis and explainability findings. Chapter~5 concludes the work and
identifies the way ahead. The appendices list the software environment, the
engineered feature set and the derived ensemble weights, and a user manual
describes how to reproduce and operate the system.

%------------------------------------------------------------------------------
\chapter{LITERATURE REVIEW/BACKGROUND STUDY}

\section{Timeline of the reported problem}

Unsolicited electronic messaging is almost as old as electronic messaging
itself: the first documented mass advertisement was sent over ARPANET in
1978. Volumes exploded with commercial e-mail in the 1990s, and the first
generation of defences was rule-based --- keyword blacklists, sender
blocklists and hand-written heuristics. The statistical era opened when
Sahami et al.\ (1998) framed junk-mail filtering as Naive-Bayes text
classification, followed quickly by margin-based methods (Drucker et al.,
1999) and, through the 2000s, by tree ensembles such as Random Forests
(Breiman, 2001). Research on the SMS channel became reproducible when
Almeida and G\'{o}mez Hidalgo (2011) released the UCI SMS Spam Collection,
while the Enron-Spam corpora (Metsis et al., 2006) did the same for
personalised e-mail filtering. From roughly 2015, deep neural models pushed
benchmark accuracy still higher but at a severe cost in transparency ---
which in turn triggered the explainable-AI wave: LIME (Ribeiro et al., 2016)
and SHAP (Lundberg and Lee, 2017) made post-hoc, per-prediction explanation
practical. Most recently, regulation (GDPR, 2018; DPDP Act, 2023) has begun
to treat explainability as a compliance requirement rather than a luxury,
which is the point in the timeline where this project positions itself.

\section{Proposed solutions}

Reported approaches fall into four broad families.

\subsection{Rule-based and collaborative filtering}

Blacklists, keyword rules and community-scored heuristic engines (e.g.\
SpamAssassin-style scoring). They are fully transparent and cheap, but
brittle: spammers evade fixed rules within days, false-positive rates are
high, and the rule bases require continuous manual curation.

\subsection{Classical statistical learning}

Naive Bayes, logistic regression, support vector machines and tree ensembles
trained on bag-of-words or TF--IDF representations. These methods reach
95--99\% accuracy on the public benchmark corpora, train in seconds to
minutes on commodity CPUs, and expose at least coarse interpretability
(coefficients, feature importances). Class imbalance --- severe on the SMS
channel --- must be handled explicitly.

\subsection{Deep learning}

Convolutional, recurrent and, more recently, transformer-based text
classifiers. They deliver the highest raw benchmark numbers but demand GPUs
and large training sets, and their decision process is opaque; attention
weights are known to be an unreliable proxy for faithful explanation.

\subsection{Explainable-AI methods}

Post-hoc techniques that explain an already-trained classifier. LIME fits a
local linear surrogate around one prediction of \emph{any} model by
perturbing the input; SHAP computes additive Shapley-value attributions, with
an exact polynomial-time algorithm (TreeSHAP) for tree ensembles. Combining a
strong classical ensemble with these two techniques promises near
state-of-the-art accuracy \emph{with} faithful explanations.

\section{Bibliometric analysis}

A comparison of the reviewed families on key features, effectiveness and
drawbacks is summarised in Table~\ref{tab:biblio}.

\begin{table}[H]
  \centering
  \caption{Comparison of reported spam-detection approaches}
  \label{tab:biblio}
  \begin{tabular}{|L{3.1cm}|L{3.4cm}|L{3.1cm}|L{4.0cm}|}
    \hline
    \textbf{Approach} & \textbf{Key feature} & \textbf{Effectiveness} & \textbf{Drawback} \\
    \hline
    Rule-based filters & Hand-written, fully transparent rules & Moderate; high false positives & Easily evaded; constant manual updating \\
    \hline
    Naive Bayes & Fast probabilistic word model & 95--98\% on benchmark corpora & Feature-independence assumption \\
    \hline
    SVM / Logistic Regression & Margin / linear decision boundary & 97--98\% & Linear view of language; needs calibration \\
    \hline
    Tree ensembles (RF, XGBoost) & Non-linear feature interactions & 97--99\% & Interpretability limited to importances \\
    \hline
    Deep / transformer models & Learned contextual representations & 98--99\%+ & Opaque; GPU and data hungry \\
    \hline
    Classical ensemble + XAI (this work) & Accuracy of ensembles with per-verdict explanations & 98.3--98.6\% (measured) & Explanation adds inference cost \\
    \hline
  \end{tabular}
\end{table}

\section{Review Summary}

The literature exhibits a persistent trade-off: the most transparent systems
are the least accurate, and the most accurate are the least transparent.
Furthermore, nearly all published work targets a single channel, even though
SMS and e-mail scams are two faces of the same problem with radically
different message statistics. A hybrid ensemble of heterogeneous classical
learners, wrapped in a dual model-agnostic/model-specific explanation layer
and specialised per modality, occupies the unclaimed middle ground: benchmark
accuracy within a point of deep models, full per-prediction transparency, and
CPU-only training. This project adopts exactly that position.

\section{Problem Definition}

To design, build and validate an interpretable multi-modal spam detection
framework that classifies both SMS and e-mail messages as legitimate or
fraudulent using a hybrid ensemble of Logistic Regression, Random Forest and
XGBoost, and that attaches a faithful local explanation (LIME over the full
ensemble; SHAP over the boosted-tree component) to every verdict, together
with global model summaries. The framework must be reproducible end-to-end on
a commodity CPU and be operable by a non-expert through an interactive
application. The work explicitly excludes non-English corpora, adversarially
obfuscated spelling, image-based spam, and hardening for production
deployment.

\section{Goals/Objectives}\label{sec:objectives}

\begin{itemize}
  \item Achieve at least 97\% test-set accuracy on both public corpora (UCI
        SMS Spam Collection; Enron-Spam).
  \item Ensure the hybrid ensemble matches or beats every base learner on
        ROC-AUC for both modalities.
  \item Attach a local explanation to every single prediction and provide
        global explanations of what each model has learned.
  \item Keep the complete pipeline --- data download, training, evaluation,
        figure generation --- reproducible on a commodity CPU in under five
        minutes, verified by an automated test suite.
  \item Deliver an interactive demonstration application usable by a
        non-expert.
\end{itemize}

%------------------------------------------------------------------------------
\chapter{DESIGN FLOW/PROCESS}

\section{Evaluation \& Selection of Specifications/Features}

The candidate features identified in the literature were critically evaluated
against the client need of \emph{accurate and auditable} filtering.

\textbf{Retained.} Modality-specific pipelines (one specialist per channel,
selected by a lightweight, user-overridable router); normalisation of
high-signal entities into \emph{readable} placeholder tokens
(\texttt{httpaddr}, \texttt{numbr}, \texttt{moneysymb}, \texttt{phonenumbr})
instead of deleting them; fifteen engineered stylometric features computed on
the raw message and named in UPPER-CASE so that they remain visually distinct
from vocabulary in every explanation plot; TF--IDF over uni- and bi-grams
with sublinear term frequency; a heterogeneous three-model ensemble
(linear~/ bagging~/ boosting) fused by performance-weighted soft voting;
balanced class weights for imbalance handling; and a dual explanation layer
(LIME + SHAP).

\textbf{Rejected.} Stemming and lemmatisation (they destroy the surface word
forms on which readable explanations depend); deletion of URLs, numbers and
currency symbols during cleaning (they are among the strongest scam signals);
deep contextual models as the primary classifier (opaque, GPU-dependent);
synthetic minority oversampling of text (it manufactures training
``messages'' that never existed and cannot be meaningfully explained); and an
on-premise mail-server integration, which added engineering cost without
addressing the stated research need.

\section{Design Constraints}

\begin{itemize}
  \item \textbf{Regulations:} transparency expectations of GDPR Article~22
        and the DPDP Act, 2023 motivate per-prediction explanations; the
        research licences of both corpora restrict the system to research and
        education use, and it is labelled accordingly.
  \item \textbf{Economic:} zero licence cost (fully open-source stack) and
        commodity-CPU training; no GPU may be assumed.
  \item \textbf{Environmental:} CPU-only training completing in about two
        minutes keeps the energy footprint negligible compared with
        GPU-trained deep models.
  \item \textbf{Health and Safety (user safety):} a false positive can
        suppress a legitimate OTP or an important mail, so high precision is
        prioritised on the SMS channel; no live user data is collected ---
        only public, anonymised research corpora are processed.
  \item \textbf{Manufacturability (reproducibility):} pinned dependencies,
        fixed random seeds, a single-command training pipeline and a
        13-check automated smoke-test suite guarantee that any examiner can
        rebuild the system.
  \item \textbf{Ethical and Social:} the scam-keyword lexicon is a plain,
        inspectable list in \texttt{src/config.py}; global SHAP summaries are
        used to audit the models for dataset bias; both corpora are credited
        to their original authors.
\end{itemize}

\section{Analysis and Feature finalization subject to constraints}

The constraint set directly shaped the final feature list. The
interpretability requirement eliminated stemming and deep contextual
encoders, and dictated the readable-placeholder normalisation: a token such
as \texttt{httpaddr} can appear verbatim inside a LIME highlight, whereas a
deleted URL cannot appear anywhere. The commodity-CPU constraint fixed the
representation at TF--IDF plus engineered statistics, giving compact feature
spaces of 3{,}015 dimensions for SMS and 5{,}015 for e-mail. The user-safety
constraint resolved the imbalance-handling question in favour of balanced
class weights and \texttt{scale\_pos\_weight} rather than synthetic
oversampling, so that every training example remains a real, explainable
message. Finally, the reproducibility constraint required every artefact ---
processed datasets, trained detectors, metrics, figures --- to be regenerated
by one command.

\section{Design Flow}

Two alternative designs were carried forward for detailed comparison.

\subsection{Design~I: Fine-tuned transformer classifier}

A single compact transformer (e.g.\ DistilBERT) fine-tuned jointly on both
corpora, with attention maps and gradient-based attributions offered as
explanations. Expected accuracy is at or marginally above the classical
state of the art, and one model serves both channels. However, fine-tuning
requires GPU time, inference is an order of magnitude slower on CPU, and ---
decisively --- attention weights are an unfaithful basis for explanation, so
the design fails the core transparency requirement of the client.

\subsection{Design~II: Hybrid classical ensemble with a dual XAI layer}

Each modality receives its own feature space and its own trained ensemble of
three heterogeneous learners --- Logistic Regression, Random Forest and
XGBoost --- combined in probability space with weights proportional to each
model's cross-validated F1. LIME explains the complete ensemble locally by
re-scoring word-dropout perturbations through the whole pipeline; exact
TreeSHAP explains the XGBoost component locally and globally. Training is
CPU-only and completes in about two minutes. The block diagram of this
alternative is shown in Figure~\ref{fig:block}.

\begin{figure}[H]
  \centering
  \begin{tikzpicture}[node distance=0.45cm and 0.7cm]
    \node[dat, minimum width=4.2cm] (raw)  {Raw message (SMS / E-mail)};
    \node[blk, below=of raw, minimum width=4.2cm] (rout) {Modality router\\(length/structure heuristic, user-overridable)};
    \node[blk, below=of rout, minimum width=7.6cm] (norm)
      {Normalisation: URLs $\to$ \texttt{httpaddr} $\cdot$ numbers $\to$ \texttt{numbr} $\cdot$ money $\to$ \texttt{moneysymb}};
    \node[blk, below=of norm, minimum width=7.6cm] (feat)
      {Feature extraction: TF--IDF (1--2 grams) + 15 engineered stylometric features};
    \node[mdl, below=1.0cm of feat] (rf) {Random\\Forest};
    \node[mdl, left=of rf]  (lr) {Logistic\\Regression};
    \node[mdl, right=of rf] (xgb) {XGBoost};
    \node[blk, below=1.0cm of rf, minimum width=6.4cm, fill=blue!12] (ens)
      {Hybrid soft-voting ensemble\\(weights $\propto$ cross-validated F1)};
    \node[dat, below left=0.75cm and -1.3cm of ens, minimum width=3.6cm, fill=green!12] (ver)
      {Verdict +\\spam probability};
    \node[xai, below right=0.75cm and -1.3cm of ens, minimum width=4.4cm] (xai)
      {XAI layer:\\LIME (full ensemble)\\SHAP (exact TreeSHAP)};
    \draw[arr] (raw) -- (rout);
    \draw[arr] (rout) -- (norm);
    \draw[arr] (norm) -- (feat);
    \draw[arr] (feat.south) -- (lr.north);
    \draw[arr] (feat.south) -- (rf.north);
    \draw[arr] (feat.south) -- (xgb.north);
    \draw[arr] (lr.south)  -- (ens.north);
    \draw[arr] (rf.south)  -- (ens.north);
    \draw[arr] (xgb.south) -- (ens.north);
    \draw[arr] (ens.south) -- (ver.north);
    \draw[arr] (ens.south) -- (xai.north);
    \draw[arr] (feat.east) .. controls +(2.2,0) and +(1.9,1.4) .. (xai.east);
  \end{tikzpicture}
  \caption{Block diagram of the proposed interpretable multi-modal detector (Design~II)}
  \label{fig:block}
\end{figure}

\section{Design selection}

The two alternatives were scored against the constraint set, as shown in
Table~\ref{tab:designsel}. Design~II was selected: it reaches the accuracy
band of the transformer on these corpora while satisfying every constraint
the transformer violates, and its tree-based component admits \emph{exact}
Shapley-value explanations rather than approximations.

\begin{table}[H]
  \centering
  \caption{Comparison and selection of the design alternatives}
  \label{tab:designsel}
  \begin{tabular}{|L{4.4cm}|L{4.3cm}|L{4.6cm}|}
    \hline
    \textbf{Criterion} & \textbf{Design~I (transformer)} & \textbf{Design~II (hybrid + XAI)} \\
    \hline
    Expected / measured accuracy & 98--99\% (reported) & 98.25--98.55\% (measured) \\
    \hline
    Explanation fidelity & Low (attention is unfaithful) & High (model-agnostic LIME; exact TreeSHAP) \\
    \hline
    Training hardware / time & GPU; tens of minutes--hours & CPU; $\approx$2 minutes \\
    \hline
    Reproducibility for examiners & Moderate & High (single command + tests) \\
    \hline
    Meets full constraint set & No & Yes \\
    \hline
  \end{tabular}
\end{table}

\section{Implementation plan/methodology}

The selected design is realised by the pipeline of Figure~\ref{fig:flow}.
Text is vectorised with sublinear TF--IDF,
\begin{equation}
  w_{t,d} \;=\; \bigl(1+\log \mathrm{tf}_{t,d}\bigr)\cdot
                \log\frac{N}{\mathrm{df}_t},
  \label{eq:tfidf}
\end{equation}
and the fifteen engineered statistics are appended to the vector. Each base
learner is trained with imbalance handling, and its 5-fold cross-validated F1
on the training split determines its voice in the ensemble:
\begin{equation}
  p(\mathrm{spam}\mid x) \;=\; \sum_{m=1}^{3} w_m\, p_m(\mathrm{spam}\mid x),
  \qquad
  w_m \;=\; \frac{F1^{\mathrm{CV}}_m}{\sum_{k=1}^{3} F1^{\mathrm{CV}}_k},
  \label{eq:ensemble}
\end{equation}
so a stronger learner gets a proportionally louder vote. The held-out test
split is touched exactly once, for the final evaluation of Chapter~4.

\begin{figure}[H]
  \centering
  \begin{tikzpicture}[node distance=0.4cm and 1.0cm]
    \node[dat, minimum width=5.4cm] (a) {Download corpora (with mirror fallbacks)};
    \node[blk, below=of a, minimum width=5.4cm] (b) {Clean, de-duplicate, label};
    \node[blk, below=of b, minimum width=5.4cm] (c) {Stratified 80/20 train--test split};
    \node[blk, below=of c, minimum width=5.4cm] (d) {Normalise text $\to$ placeholder tokens};
    \node[blk, right=of a, minimum width=5.9cm] (e) {Featurise: TF--IDF + 15 stylometric};
    \node[blk, below=of e, minimum width=5.9cm] (f) {5-fold CV per base learner $\Rightarrow$ weights $w_m$};
    \node[blk, below=of f, minimum width=5.9cm] (g) {Fit LR, RF, XGBoost; build weighted ensemble};
    \node[blk, below=of g, minimum width=5.9cm] (h) {Evaluate once on held-out test set};
    \node[xai, below=of h, minimum width=5.9cm] (i) {Generate SHAP/LIME artefacts, figures, metrics};
    \draw[arr] (a) -- (b);
    \draw[arr] (b) -- (c);
    \draw[arr] (c) -- (d);
    \draw[arr] (d.east) .. controls +(0.9,0) and +(-0.9,0) .. (e.west);
    \draw[arr] (e) -- (f);
    \draw[arr] (f) -- (g);
    \draw[arr] (g) -- (h);
    \draw[arr] (h) -- (i);
  \end{tikzpicture}
  \caption{Implementation flowchart of the training and evaluation pipeline}
  \label{fig:flow}
\end{figure}

%------------------------------------------------------------------------------
\chapter{RESULTS ANALYSIS AND VALIDATION}

\section{Implementation of solution}

Modern engineering tools were used throughout:

\begin{itemize}
  \item \textbf{Analysis:} Python (pandas, NumPy) and a Jupyter research
        notebook for exploratory data analysis and sampling studies.
  \item \textbf{Design and implementation:} scikit-learn (Logistic
        Regression, Random Forest, TF--IDF), XGBoost, and the SHAP and LIME
        libraries, organised as a modular \texttt{src/} package.
  \item \textbf{Report preparation:} \LaTeX{} for this document; matplotlib
        with a shared publication style for all generated figures.
  \item \textbf{Project management and communication:} Git for version
        control and a shared Kanban board for task tracking.
  \item \textbf{Testing and validation:} a 13-check automated end-to-end
        smoke-test suite (including a simulated application session) and a
        strict held-out evaluation protocol.
\end{itemize}

\subsection{Experimental setup}

Two public research corpora were used; their statistics after cleaning and
de-duplication are given in Table~\ref{tab:datasets} and visualised in
Figure~\ref{fig:classbalance}. Each corpus received an 80/20 stratified
train--test split; base-learner weights were computed by 5-fold
cross-validation on the training split only, and the test split was used
exactly once. Performance is reported with accuracy, precision
$P=\mathrm{TP}/(\mathrm{TP}+\mathrm{FP})$, recall
$R=\mathrm{TP}/(\mathrm{TP}+\mathrm{FN})$, their harmonic mean F1, and
ROC-AUC.

\begin{table}[H]
  \centering
  \caption{Corpora used for training and evaluation (after cleaning)}
  \label{tab:datasets}
  \begin{tabular}{|L{1.7cm}|L{5.6cm}|c|c|c|}
    \hline
    \textbf{Modality} & \textbf{Corpus} & \textbf{Messages} & \textbf{Spam share} & \textbf{Median words} \\
    \hline
    SMS & UCI SMS Spam Collection (Almeida \& G\'{o}mez Hidalgo, 2011) & 5,158 & 12.4\% & 12 \\
    \hline
    E-mail & Enron-Spam, stratified subsample (Metsis et al., 2006) & 12,000 & 48.3\% & 155 \\
    \hline
  \end{tabular}
\end{table}

\begin{figure}[H]
  \centering
  \figimg[0.9\textwidth]{figures/dataset_overview.png}{2.4in}
  \caption{Class balance per modality and message-length distributions}
  \label{fig:classbalance}
\end{figure}

The two modalities differ radically in both imbalance (12\% versus 48\%
spam) and length (12 versus 155 median words), which is precisely
why each is served by its own feature space and its own trained ensemble.

\subsection{Classification results}

Tables~\ref{tab:sms} and \ref{tab:email} report the complete held-out
results for all four systems on each modality; Figure~\ref{fig:metrics}
compares them graphically and Figure~\ref{fig:roc} shows the ROC curves.

\begin{table}[H]
  \centering
  \caption{SMS --- held-out test set (1,032 messages)}
  \label{tab:sms}
  \setlength{\tabcolsep}{4pt}
  \begin{tabular}{|L{3.3cm}|c|c|c|c|c|c|}
    \hline
    \textbf{Model} & \textbf{Accuracy} & \textbf{Precision} & \textbf{Recall} & \textbf{F1} & \textbf{ROC-AUC} & \textbf{CV F1} \\
    \hline
    Logistic Regression & 0.9855 & 0.9520 & 0.9297 & 0.9407 & 0.9941 & 0.9554 $\pm$ 0.020 \\
    \hline
    Random Forest & 0.9845 & 0.9516 & 0.9219 & 0.9365 & 0.9929 & 0.9558 $\pm$ 0.017 \\
    \hline
    XGBoost & 0.9845 & 0.9667 & 0.9062 & 0.9355 & 0.9941 & 0.9586 $\pm$ 0.015 \\
    \hline
    \textbf{Hybrid Ensemble} & \textbf{0.9855} & \textbf{0.9669} & \textbf{0.9141} & \textbf{0.9398} & \textbf{0.9941} & --- \\
    \hline
  \end{tabular}
\end{table}

\begin{table}[H]
  \centering
  \caption{E-mail --- held-out test set (2,400 messages)}
  \label{tab:email}
  \setlength{\tabcolsep}{4pt}
  \begin{tabular}{|L{3.3cm}|c|c|c|c|c|c|}
    \hline
    \textbf{Model} & \textbf{Accuracy} & \textbf{Precision} & \textbf{Recall} & \textbf{F1} & \textbf{ROC-AUC} & \textbf{CV F1} \\
    \hline
    Logistic Regression & 0.9833 & 0.9753 & 0.9905 & 0.9829 & 0.9979 & 0.9797 $\pm$ 0.004 \\
    \hline
    Random Forest & 0.9746 & 0.9629 & 0.9853 & 0.9740 & 0.9973 & 0.9669 $\pm$ 0.004 \\
    \hline
    XGBoost & 0.9750 & 0.9645 & 0.9845 & 0.9744 & 0.9969 & 0.9732 $\pm$ 0.003 \\
    \hline
    \textbf{Hybrid Ensemble} & \textbf{0.9825} & \textbf{0.9705} & \textbf{0.9940} & \textbf{0.9821} & \textbf{0.9983} & --- \\
    \hline
  \end{tabular}
\end{table}

\begin{figure}[H]
  \centering
  \begin{tikzpicture}
    \begin{axis}[
      ybar, bar width=4.2pt, width=0.44\textwidth, height=5.6cm,
      ymin=0.90, ymax=1.004,
      symbolic x coords={LR,RF,XGB,Ensemble}, xtick=data,
      enlarge x limits=0.16, ymajorgrids,
      title={SMS (test, $n=1032$)},
      tick label style={font=\scriptsize},
      label style={font=\scriptsize},
      title style={font=\footnotesize\bfseries},
      ylabel={Score},
      legend style={font=\scriptsize, at={(1.09,-0.24)}, anchor=north, legend columns=3},
    ]
      \addplot coordinates {(LR,0.9855)(RF,0.9845)(XGB,0.9845)(Ensemble,0.9855)};
      \addplot coordinates {(LR,0.9407)(RF,0.9365)(XGB,0.9355)(Ensemble,0.9398)};
      \addplot coordinates {(LR,0.9941)(RF,0.9929)(XGB,0.9941)(Ensemble,0.9941)};
      \legend{Accuracy,F1,ROC-AUC}
    \end{axis}
  \end{tikzpicture}\hfill
  \begin{tikzpicture}
    \begin{axis}[
      ybar, bar width=4.2pt, width=0.44\textwidth, height=5.6cm,
      ymin=0.90, ymax=1.004,
      symbolic x coords={LR,RF,XGB,Ensemble}, xtick=data,
      enlarge x limits=0.16, ymajorgrids,
      title={E-mail (test, $n=2400$)},
      tick label style={font=\scriptsize},
      label style={font=\scriptsize},
      title style={font=\footnotesize\bfseries},
    ]
      \addplot coordinates {(LR,0.9833)(RF,0.9746)(XGB,0.9750)(Ensemble,0.9825)};
      \addplot coordinates {(LR,0.9829)(RF,0.9740)(XGB,0.9744)(Ensemble,0.9821)};
      \addplot coordinates {(LR,0.9979)(RF,0.9973)(XGB,0.9969)(Ensemble,0.9983)};
    \end{axis}
  \end{tikzpicture}
  \caption{Test-set metrics per model and modality (vertical axis truncated at 0.90 to resolve the differences)}
  \label{fig:metrics}
\end{figure}

\begin{figure}[H]
  \centering
  \figimg[0.48\textwidth]{figures/sms_roc_curves.png}{2.3in}\hfill
  \figimg[0.48\textwidth]{figures/email_roc_curves.png}{2.3in}
  \caption{ROC curves for all four systems: SMS (left) and e-mail (right)}
  \label{fig:roc}
\end{figure}

On ROC-AUC the hybrid ensemble ties the best base learners on SMS (0.9941,
identical to four decimal places) and wins outright on e-mail (0.9983), where
it also posts near-perfect recall (0.9940) at 0.9705 precision. On SMS it
pairs the highest precision of any system (0.9669) with the joint-highest
accuracy (0.9855, tied with Logistic Regression) --- exactly the profile the
user-safety constraint asked for, since on that channel a false alarm means a
suppressed legitimate message.

\subsection{Ensemble behaviour and error analysis}

Figure~\ref{fig:confusion} shows the ensemble confusion matrices derived from
the held-out evaluations. The SMS ensemble mislabels only 4 of 904 legitimate
messages; the e-mail ensemble misses just 7 of 1,158 spam messages. The two
matrices are deliberately asymmetric: the SMS specialist errs on the side of
precision, the e-mail specialist on the side of recall, matching the relative
costs of a wrong decision on each channel.

\begin{figure}[H]
  \centering
  \begin{minipage}{0.48\textwidth}\centering
  \begin{tikzpicture}[font=\footnotesize]
    \def\cw{2.15}\def\ch{1.2}
    \fill[green!14] (0,\ch) rectangle (\cw,2*\ch);
    \fill[red!12]   (\cw,\ch) rectangle (2*\cw,2*\ch);
    \fill[red!12]   (0,0) rectangle (\cw,\ch);
    \fill[green!14] (\cw,0) rectangle (2*\cw,\ch);
    \draw (0,0) rectangle (2*\cw,2*\ch);
    \draw (\cw,0) -- (\cw,2*\ch);
    \draw (0,\ch) -- (2*\cw,\ch);
    \node[align=center] at (0.5*\cw,1.5*\ch) {TN\\\textbf{900}};
    \node[align=center] at (1.5*\cw,1.5*\ch) {FP\\\textbf{4}};
    \node[align=center] at (0.5*\cw,0.5*\ch) {FN\\\textbf{11}};
    \node[align=center] at (1.5*\cw,0.5*\ch) {TP\\\textbf{117}};
    \node[anchor=south] at (0.5*\cw,2*\ch+0.05) {Pred.\ ham};
    \node[anchor=south] at (1.5*\cw,2*\ch+0.05) {Pred.\ spam};
    \node[anchor=east]  at (-0.06,1.5*\ch) {Actual ham};
    \node[anchor=east]  at (-0.06,0.5*\ch) {Actual spam};
    \node[anchor=south, font=\footnotesize\bfseries] at (\cw,2*\ch+0.5) {SMS ensemble ($n=1032$)};
  \end{tikzpicture}
  \end{minipage}\hfill
  \begin{minipage}{0.48\textwidth}\centering
  \begin{tikzpicture}[font=\footnotesize]
    \def\cw{2.15}\def\ch{1.2}
    \fill[green!14] (0,\ch) rectangle (\cw,2*\ch);
    \fill[red!12]   (\cw,\ch) rectangle (2*\cw,2*\ch);
    \fill[red!12]   (0,0) rectangle (\cw,\ch);
    \fill[green!14] (\cw,0) rectangle (2*\cw,\ch);
    \draw (0,0) rectangle (2*\cw,2*\ch);
    \draw (\cw,0) -- (\cw,2*\ch);
    \draw (0,\ch) -- (2*\cw,\ch);
    \node[align=center] at (0.5*\cw,1.5*\ch) {TN\\\textbf{1207}};
    \node[align=center] at (1.5*\cw,1.5*\ch) {FP\\\textbf{35}};
    \node[align=center] at (0.5*\cw,0.5*\ch) {FN\\\textbf{7}};
    \node[align=center] at (1.5*\cw,0.5*\ch) {TP\\\textbf{1151}};
    \node[anchor=south] at (0.5*\cw,2*\ch+0.05) {Pred.\ ham};
    \node[anchor=south] at (1.5*\cw,2*\ch+0.05) {Pred.\ spam};
    \node[anchor=east]  at (-0.06,1.5*\ch) {Actual ham};
    \node[anchor=east]  at (-0.06,0.5*\ch) {Actual spam};
    \node[anchor=south, font=\footnotesize\bfseries] at (\cw,2*\ch+0.5) {E-mail ensemble ($n=2400$)};
  \end{tikzpicture}
  \end{minipage}
  \caption{Confusion matrices of the hybrid ensembles on the held-out test sets}
  \label{fig:confusion}
\end{figure}

\subsection{Explainability results}

\subsubsection{A verdict, explained (local scope)}

The following scam SMS, never seen in training, is used as a live exhibit:

\begin{center}
\fbox{\parbox{0.88\textwidth}{\small
``URGENT! You have WON a \pounds 2,000 cash prize GUARANTEED. Call
09061701461 now or visit www.claim-your-prize.com to claim!''}}
\end{center}

All three base learners vote spam --- Logistic Regression 1.000, Random
Forest 0.987, XGBoost 1.000 --- giving a weighted ensemble probability of
\textbf{0.996 spam}. LIME re-scores approximately 1,000 word-dropout
perturbations of the message through the entire pipeline and fits a local
linear surrogate; the resulting per-word contributions are plotted in
Figure~\ref{fig:lime} alongside the SHAP waterfall of the same prediction.
The premium-rate phone number is the single strongest signal, followed by the
money amount and the classic scam vocabulary --- a justification a non-expert
can check at a glance.

\begin{figure}[H]
  \centering
  \begin{minipage}[c]{0.5\textwidth}\centering
  \begin{tikzpicture}
    \begin{axis}[
      xbar, width=0.76\textwidth, height=6.4cm,
      xmin=0, xmax=0.15,
      symbolic y coords={your,www,WON,GUARANTEED,2,prize,URGENT,claim,000,09061701461},
      ytick=data,
      yticklabel style={font=\scriptsize\ttfamily},
      tick label style={font=\scriptsize},
      xlabel={LIME weight toward ``spam''},
      label style={font=\scriptsize},
      nodes near coords, nodes near coords style={font=\tiny},
      /pgf/number format/.cd, fixed, precision=3,
      bar width=5.5pt,
    ]
      \addplot coordinates {
        (0.025,your) (0.031,www) (0.036,WON) (0.037,GUARANTEED) (0.038,2)
        (0.045,prize) (0.046,URGENT) (0.053,claim) (0.061,000) (0.123,09061701461)};
    \end{axis}
  \end{tikzpicture}
  \end{minipage}\hfill
  \begin{minipage}[c]{0.46\textwidth}\centering
  \figimg[0.98\textwidth]{figures/sms_shap_waterfall.png}{2.4in}
  \end{minipage}
  \caption{Local explanations of the exhibited SMS: LIME word attributions over the full ensemble (left) and SHAP waterfall of the XGBoost component (right)}
  \label{fig:lime}
\end{figure}

\subsubsection{What the models learned (global scope)}

Global SHAP summaries were computed over 800 test messages per modality
(Figure~\ref{fig:beeswarm}). The key finding is that \emph{stylometry rivals
vocabulary}: on SMS the five top-ranked features are all engineered
stylometric statistics --- \texttt{DIGIT\_RATIO}, \texttt{AVG\_WORD\_LEN},
\texttt{SPAM\_KEYWORDS}, \texttt{UPPERCASE\_RATIO} and \texttt{CHAR\_COUNT}
--- and the fifteen engineered features together carry about three-quarters
(74.5\%) of the total attribution, although they are only 15 of the 3,015
features. A scam is betrayed by its \emph{shape} as much as by its words.

All SHAP values in this chapter are computed with XGBoost's native, exact
TreeSHAP on the same sparse feature matrix that the model scores. This
matters: XGBoost treats entries that are absent from a sparse matrix as
missing values, so explaining a densified copy of the input would describe a
different function --- in this pipeline, densification changes the XGBoost
verdict for 729 of the 1,032 SMS test messages.

\begin{figure}[H]
  \centering
  \figimg[0.48\textwidth]{figures/sms_shap_beeswarm.png}{2.6in}\hfill
  \figimg[0.48\textwidth]{figures/email_shap_beeswarm.png}{2.6in}
  \caption{Global SHAP beeswarm summaries: SMS (left) and e-mail (right); UPPER-CASE rows are engineered stylometric features, lower-case rows are words}
  \label{fig:beeswarm}
\end{figure}

\subsubsection{XAI as a debugging tool: a dataset bias, caught}

On the e-mail model, the global SHAP ranking exposed the token
\texttt{enron} as a strong \emph{ham} indicator: the corpus's legitimate mail
is all Enron-internal, so the model partly learned the dataset rather than
the concept. This bias is invisible in every accuracy table in this chapter
and obvious in one beeswarm plot --- the clearest argument this project makes
for shipping explanations together with predictions.

\subsection{Interpretation and data validation}

Cross-validation dispersion is small (F1 standard deviation between
$\pm$0.003 and $\pm$0.020 across all six base learners), indicating stable
models rather than lucky splits, and the derived ensemble weights are
near-uniform (Appendix-3), so the ensemble's gains come from the
\emph{diversity} of the three learners rather than from one dominant model.
The 13-check automated smoke-test suite exercises the entire system ---
data loading, preprocessing, both detectors, LIME, SHAP and a simulated
application session --- and passes end-to-end; the complete pipeline
reproduces from a clean checkout in about three minutes on a commodity CPU,
with training itself taking about two minutes. The dataset downloader
validates its sources and falls back to mirrors automatically, and the
modality router can be overridden by the user in the application, so a
misrouted message can always be re-scored by the other specialist.

%------------------------------------------------------------------------------
\chapter{CONCLUSION AND FUTURE WORK}
\vspace{-12pt}   % keeps the last future-work bullet off a page of its own

\section{Conclusion}

Every objective set in Section~\ref{sec:objectives} was met. The hybrid ensemble reached
98.55\% accuracy on SMS and 98.25\% on e-mail against a target of 97\%, and
it matched (to four decimal places) or beat every base learner on ROC-AUC on
both modalities (0.9941
on SMS; an outright best 0.9983 on e-mail). Every single prediction is
delivered with a local LIME explanation over the full ensemble and a SHAP
attribution of the boosted-tree component, and global beeswarm and importance
summaries describe what each model has learned. The complete pipeline ---
corpus download, training, evaluation and figure generation --- reproduces
from a clean checkout in about three minutes on a commodity CPU and is
verified by a 13-check automated test suite, and the four-tab Streamlit
application makes the system operable by a non-expert.

The most instructive result, however, was not a number. The global SHAP
analysis exposed the token \texttt{enron} acting as a shortcut for legitimate
mail --- a genuine dataset bias in the e-mail corpus that deviated from the
naive expectation that a 98\%-accurate model has learned the concept of
``spam''. The bias is a property of the data, not a fault of the pipeline,
but it would have shipped silently inside any black-box deployment. Catching
it required the explanation layer, which is precisely the thesis this
project set out to demonstrate: interpretability is not a presentation
feature bolted onto a classifier, it is a debugging and auditing instrument
that changes what the engineer can see.

\section{Future work}

\begin{itemize}[itemsep=0pt]
  \item Extend the framework to multilingual and code-mixed corpora, since
        the present datasets are English-only and date from 2002--2012.
  \item Perform cross-domain evaluation (train on one corpus, test on
        another) to quantify and then mitigate the dataset bias exposed by
        SHAP.
  \item Add probability calibration so that the reported spam probabilities
        can be thresholded meaningfully by downstream policy.
  \item Harden the models against adversarially obfuscated spelling (e.g.\
        leet-speak substitutions) and measure robustness explicitly.
  \item Add counterfactual explanations --- ``what is the smallest edit that
        flips this verdict?'' --- as a third, complementary XAI technique.
  \item Benchmark a fine-tuned transformer under the same explanation
        harness to quantify the accuracy--interpretability trade-off
        empirically.
\end{itemize}

%==============================================================================
%  REFERENCES
%==============================================================================
\chapter*{REFERENCES}
\phantomsection\addcontentsline{toc}{chapter}{REFERENCES}

\begin{singlespace}
\begin{enumerate}[leftmargin=1.4cm]
  \item Almeida, T.A. and G\'{o}mez Hidalgo, J.M. (2011) `SMS Spam Collection v.1', UCI Machine Learning Repository.
  \item Breiman, L. (2001) `Random Forests', Machine Learning, Vol.~45, No.~1, pp.~5--32.
  \item Chen, T. and Guestrin, C. (2016) `XGBoost: A Scalable Tree Boosting System', Proc.~22nd ACM SIGKDD International Conference on Knowledge Discovery and Data Mining, pp.~785--794.
  \item Cranor, L.F. and LaMacchia, B.A. (1998) `Spam!', Communications of the ACM, Vol.~41, No.~8, pp.~74--83.
  \item Drucker, H., Wu, D. and Vapnik, V.N. (1999) `Support Vector Machines for Spam Categorization', IEEE Transactions on Neural Networks, Vol.~10, No.~5, pp.~1048--1054.
  \item Federal Bureau of Investigation (2024) `Internet Crime Report 2023', Internet Crime Complaint Center (IC3), Washington, DC.
  \item Lundberg, S.M. and Lee, S.-I. (2017) `A Unified Approach to Interpreting Model Predictions', Advances in Neural Information Processing Systems 30, pp.~4765--4774.
  \item Metsis, V., Androutsopoulos, I. and Paliouras, G. (2006) `Spam Filtering with Naive Bayes --- Which Naive Bayes?', Proc.~3rd Conference on Email and Anti-Spam (CEAS), Mountain View, CA.
  \item Pedregosa, F. et al. (2011) `Scikit-learn: Machine Learning in Python', Journal of Machine Learning Research, Vol.~12, pp.~2825--2830.
  \item Ribeiro, M.T., Singh, S. and Guestrin, C. (2016) `\,``Why Should I Trust You?'': Explaining the Predictions of Any Classifier', Proc.~22nd ACM SIGKDD International Conference on Knowledge Discovery and Data Mining, pp.~1135--1144.
  \item Sahami, M., Dumais, S., Heckerman, D. and Horvitz, E. (1998) `A Bayesian Approach to Filtering Junk E-mail', AAAI Workshop on Learning for Text Categorization, Madison, WI.
  \item Telecom Regulatory Authority of India (2018) `Telecom Commercial Communications Customer Preference Regulations', New Delhi.
\end{enumerate}
\end{singlespace}

%==============================================================================
%  APPENDIX
%==============================================================================
\chapter*{APPENDIX}
\phantomsection\addcontentsline{toc}{chapter}{APPENDIX}

\begin{center}\textbf{\fourteen Appendix-1: Software environment}\end{center}

\vspace{6pt}
\begin{singlespace}
\begin{center}
\begin{tabular}{|c|L{4.0cm}|L{7.6cm}|}
  \hline
  \textbf{S.~No.} & \textbf{Component} & \textbf{Role in the framework} \\
  \hline
  1 & Python 3.10--3.12 & Implementation language of the entire pipeline \\
  \hline
  2 & scikit-learn & Logistic Regression, Random Forest, TF--IDF, metrics, cross-validation \\
  \hline
  3 & XGBoost & Gradient-boosted tree base learner \\
  \hline
  4 & SHAP & Exact TreeSHAP local and global attributions \\
  \hline
  5 & LIME & Model-agnostic local explanations of the full ensemble \\
  \hline
  6 & pandas / NumPy & Data handling and numerical computation \\
  \hline
  7 & matplotlib & All generated figures (shared publication style) \\
  \hline
  8 & Streamlit & Four-tab interactive demonstration application \\
  \hline
  9 & Git & Version control and collaboration \\
  \hline
\end{tabular}
\end{center}
\end{singlespace}

\vspace{14pt}
\begin{center}\textbf{\fourteen Appendix-2: Engineered stylometric features (representative)}\end{center}

\vspace{6pt}
\begin{singlespace}
\begin{center}
\begin{tabular}{|c|L{4.6cm}|L{7.0cm}|}
  \hline
  \textbf{S.~No.} & \textbf{Feature} & \textbf{Description} \\
  \hline
  1 & \texttt{DIGIT\_RATIO} & Fraction of characters that are digits \\
  \hline
  2 & \texttt{UPPERCASE\_RATIO} & Fraction of characters in upper case \\
  \hline
  3 & \texttt{CHAR\_COUNT} & Message length in characters \\
  \hline
  4 & \texttt{AVG\_WORD\_LEN} & Average word length \\
  \hline
  5 & \texttt{NUM\_EXCLAMATIONS} & Count of exclamation marks \\
  \hline
  6 & \texttt{NUM\_URLS} & Number of URLs detected \\
  \hline
  7 & \texttt{NUM\_PHONES} & Number of phone numbers detected \\
  \hline
  8 & \texttt{SPAM\_KEYWORDS} & Hits against the transparent scam-keyword lexicon \\
  \hline
\end{tabular}
\end{center}
\end{singlespace}

The full set of fifteen features, together with the complete keyword lexicon,
is defined in plain sight in \texttt{src/config.py} and
\texttt{src/features.py} so that it can be audited without reading model
internals.

\vspace{14pt}
\begin{center}\textbf{\fourteen Appendix-3: Derived ensemble weights (normalised CV F1)}\end{center}

\vspace{6pt}
\begin{singlespace}
\begin{center}
\begin{tabular}{|L{3.6cm}|c|c|c|c|}
  \hline
  \textbf{Modality} & \textbf{Base learner} & \textbf{CV F1} & \textbf{Weight $w_m$} & \\
  \cline{1-4}
  \multirow{3}{*}{SMS} & Logistic Regression & 0.9554 & 0.3329 & \multirow{3}{*}{$\sum w_m = 1$} \\
  \cline{2-4}
   & Random Forest & 0.9558 & 0.3331 & \\
  \cline{2-4}
   & XGBoost & 0.9586 & 0.3340 & \\
  \hline
  \multirow{3}{*}{E-mail} & Logistic Regression & 0.9797 & 0.3356 & \multirow{3}{*}{$\sum w_m = 1$} \\
  \cline{2-4}
   & Random Forest & 0.9669 & 0.3311 & \\
  \cline{2-4}
   & XGBoost & 0.9732 & 0.3333 & \\
  \hline
\end{tabular}
\end{center}
\end{singlespace}

The near-uniform weights show that the three learners are comparably strong;
the ensemble's advantage therefore comes from their diversity (linear~/
bagging~/ boosting) rather than from any single dominant model.

%==============================================================================
%  USER MANUAL
%==============================================================================
\chapter*{USER MANUAL}
\phantomsection\addcontentsline{toc}{chapter}{USER MANUAL}

\textbf{Prerequisites.} Python 3.10--3.12, about 2~GB of free disk space, and
an internet connection for the first run (the two public corpora are
downloaded automatically, with mirror fallbacks).

\begin{enumerate}[leftmargin=1.4cm]
  \item Obtain the project source code and open a terminal in the project
        root directory.
  \item Create and activate a virtual environment:\\
        \texttt{python -m venv .venv \&\& source .venv/bin/activate}\\
        (on Windows: \texttt{.venv\textbackslash Scripts\textbackslash activate}).
  \item Install the pinned dependencies:\\
        \texttt{pip install -r requirements.txt}
  \item Build the datasets and train both detectors (about two minutes;
        downloads the corpora on the first run):\\
        \texttt{python -m src.data\_loader}\\
        \texttt{python -m src.train}
  \item Verify the installation end-to-end (13 automated checks, including a
        simulated application session):\\
        \texttt{python tests/smoke\_test.py}
  \item Launch the demonstration application:\\
        \texttt{streamlit run app.py} \quad then open
        \texttt{http://localhost:8501} in a browser.
  \item In the \textbf{Detector} tab, paste any message (or load a real
        held-out example) and press \emph{Analyse}. The verdict card shows
        the ensemble spam probability, each base model's vote with its
        weight, the message with LIME-highlighted words, the SHAP waterfall
        and the engineered feature values (Figure~\ref{fig:dash}).
  \item Use the \textbf{Model performance} tab for the metric tables, ROC
        curves and confusion matrices; the \textbf{Global explanations} tab
        for the interactive beeswarm and importance analyses; and the
        \textbf{About} tab for the architecture and references.
  \item If the modality router picks the wrong specialist for an unusual
        message, override it manually with the SMS/E-mail selector and
        re-analyse. To rebuild everything from scratch, run
        \texttt{python -m src.data\_loader --force} and repeat step~4.
\end{enumerate}

\begin{figure}[H]
  \centering
  \figimg[0.85\textwidth]{figures/app_detector_tab.png}{2.6in}
  \caption{Detector tab of the application after analysing a message}
  \label{fig:dash}
\end{figure}

\end{document}